\textbf{Rebalancing.} Orders are placed once a month at the closing auction of the last trading day, using the weights computed
from that day's closing prices. All seven instruments (six ETFs and a T-bill ETF) trade billions of dollars a day, so a
portfolio of any realistic private-client size can be traded without moving prices. A model portfolio can be applied to
many client accounts at once.

\textbf{Turnover and costs.} The strategy trades \res{strat_turnover} of its value per year, against \res{b6040_turnover}
for 60/40: risk-parity weights change with volatility, and trend signals switch whole positions in and out. At our assumed
\res{cost_bps} per unit traded this costs \res{cost_drag} p.a.; the gross Sharpe ratio of \res{strat_gross_sharpe} falls to
\res{strat_sharpe} net. Solving $SR_{\text{strategy}}(c)=SR_{60/40}(c)$ with \texttt{fsolve} gives a \emph{break-even cost} of
\res{breakeven_bps}: above it, the 60/40 portfolio offers the better risk-adjusted return. Bid-ask spreads of these ETFs are
typically one basis point or less and commissions are often zero, so a cost-efficient platform is essential. At 50 bps the
Sharpe ratio would drop to \res{sharpe_cost50}. Simple ways to cut turnover are trading only weight changes above a threshold
and switching T-bill positions within a money-market sweep.

\textbf{Euro-based clients.} US-listed ETFs are not available to EU retail investors, so the UCITS equivalents in
Table~\ref{tab:holdings} are used. The backtest is in US dollars; a euro investor bears the EUR/USD exchange rate on the
whole portfolio. Because the strategy's own volatility is low, currency risk would dominate, so the bond and cash sleeves
should be held in EUR-hedged share classes. US investors should note that DBC issues a Schedule K-1 tax form; the
no-K-1 version (PDBC) tracks a similar index.
